\section{Tubular coordinates around the toroidal electron}
\label{app:tubular-electron}

The Williamson--van der Mark picture of the electron --- a
Compton-frequency internal mode confined to a narrow strip of
electromagnetic energy that is bent and twisted into a toroidal hull
--- can be made mathematically explicit by introducing tubular
coordinates around a closed centerline $\mathbf{C}(z)$. The physical
structure is this \emph{twisted strip}, a Möbius-type strip in the sense that the internal polarization frame is single-valued only after two circuits, which winds twice around the
torus core before closing and therefore requires a $4\pi$ rotation
(two full circuits) to return to its starting configuration. For
concrete calculations we thicken this strip slightly into a tubular
neighbourhood and parametrize it with coordinates $(z,x,y)$, where
$z$ is arclength along the centerline and $(x,y)$ span the
cross-section. It is important to emphasize that ``tube'' or ``tubular
coordinates'' refer solely to this \emph{mathematical coordinate
patch}, introduced to handle curvature and torsion in a systematic
way. The physical support of the internal mode---the region where
$|\xi|$ is appreciable---is the narrow twisted strip itself, not a
separate material tube or tunnel with walls.

Let $\mathbf{C}(z)\in\mathbb R^3$ denote the spatial centerline of the
toroidal electron in the brane's three lateral directions
$X^1,X^2,X^3$, parametrized by arclength $z\in[0,L)$ in the electron
rest frame, so that $|\partial_z\mathbf C(z)|=1$ and $\mathbf C(z+L)
=\mathbf C(z)$. This closed curve is arranged so that it winds twice
around the torus core in a $(2,1)$ torus-knot pattern: completing one
circuit of length $L$ corresponds to two revolutions around the core,
which is the geometric origin of the $4\pi$ periodicity of the internal
field configuration discussed in the main text. Along this closed curve
we introduce the Frenet--Serret frame
$(\mathbf t(z),\mathbf n(z),\mathbf b(z))$ with curvature $\kappa(z)$
and torsion $\tau(z)$,
%
\begin{equation}
  \mathbf t'=\kappa\,\mathbf n,\qquad
  \mathbf n'=-\kappa\,\mathbf t+\tau\,\mathbf b,\qquad
  \mathbf b'=-\tau\,\mathbf n,
  \label{eq:frenet-serret}
\end{equation}
%
where $'=\partial_z$. A small tubular neighbourhood of the centerline
is then parametrized by coordinates $(z,x,y)$ via
%
\begin{equation}
  \mathbf r(z,x,y)
  \;=\;
  \mathbf C(z) + x\,\mathbf n(z) + y\,\mathbf b(z),
  \label{eq:tubular-embedding}
\end{equation}
%
with $(x,y)$ spanning the cross-section of the ``tunnel'' around the
electron. The microscopic embedding field is still
$\mathbf X=(X^1,X^2,X^3,X^4)$, but in these adapted coordinates the
three lateral components $(X^1,X^2,X^3)$ are replaced by
$\mathbf r(z,x,y)$, while the transverse embedding direction (represented by $X^4$ in this gauge) remains a
separate normal coordinate. We denote the microscopic transverse-displacement mode in
this tubular neighbourhood by
%
\begin{equation}
  \xi(t,z,x,y) \;:=\; X^4\bigl(t,\mathbf r(z,x,y)\bigr),
  \label{eq:xi-from-X4}
\end{equation}
%
so that $(t,z,x,y)$ serve as local coordinates on the world-tube of
the electron---that is, the spacetime region of the brane where the
amplitude $|\xi|$ is appreciable---with the double-loop structure
encoded in the internal mode pattern of $\xi$ across the cross-section
$(x,y)$.

The induced metric $g_{ab}$ on the brane in the coordinates
$(z,x,y)$ is obtained by pulling back the flat ambient metric
in the lateral directions:
%
\begin{equation}
  g_{ab}(z,x,y)
  \;=\;
  \partial_a\mathbf r\cdot\partial_b\mathbf r,
  \qquad a,b\in\{z,x,y\}.
  \label{eq:tubular-metric-def}
\end{equation}
%
Using the Frenet--Serret relations~\eqref{eq:frenet-serret} one finds
to first order in the transverse distance $\rho=\sqrt{x^2+y^2}$ from
the centerline,
%
\begin{equation}
  g_{ab}
  \;\approx\;
  \begin{pmatrix}
    1 - 2\kappa x & -\,\tau y & \;\;\tau x \\[0.4ex]
    -\,\tau y & 1 & 0 \\[0.4ex]
    \;\;\tau x & 0 & 1
  \end{pmatrix},
  \qquad
  g^{ab}
  \;\approx\;
  \begin{pmatrix}
    1 + 2\kappa x & \;\;\tau y & -\,\tau x \\[0.4ex]
    \;\;\tau y & 1 & 0 \\[0.4ex]
    -\,\tau x & 0 & 1
  \end{pmatrix},
  \label{eq:tubular-metric}
\end{equation}
%
where $g^{ab}$ is the inverse metric, and we have neglected terms of
order $\kappa^2\rho^2$ and $\tau^2\rho^2$. In these adapted
coordinates the toroidal soliton appears, locally, as a straight tunnel
aligned with the $z$-axis; the global curvature and double-loop
structure are encoded entirely in the $z$-dependence of the frame
$(\mathbf n,\mathbf b)$ and the metric coefficients $g_{ab}(z,x,y)$.
The electromagnetic energy circulating along the double loop is then
represented by a structured internal mode $\xi(t,z,x,y)$ localized
within this world-tube (the brane region where $|\xi|$ is large),
confined to a finite range of $(x,y)$ around the centerline.

\paragraph{Effective amplitude dynamics in tubular coordinates.}

Restricting attention to the dominant microscopic amplitude mode
$\xi(t,z,x,y)$ inside the electron's world-tube, we can write a
minimal effective Lagrangian density in the adapted coordinates
$(t,z,x,y)$ of the form
%
\begin{equation}
  \mathcal{L}_\xi
  \;=\;
  \frac{\rho_m}{2}\,(\partial_t \xi)^2
  \;-\;
  \frac{T_\mathrm{eff}}{2}\,g^{ab}\,\partial_a\xi\,\partial_b\xi
  \;-\;
  V_\mathrm{nl}(\xi,\partial\xi,\dots),
  \label{eq:xi-tubular-lagrangian}
\end{equation}
%
where $\rho_m$ is the effective mass density introduced in
Section~\ref{sec:emergent-relativity}, $T_\mathrm{eff}$ is an effective
tension along the tube that depends on the microscopic spring constant
and rest length, and $V_\mathrm{nl}$ collects residual nonlinear
corrections beyond the geometric contributions contained in the metric
$g^{ab}$. Inserting the approximate inverse metric
\eqref{eq:tubular-metric} yields, to first order in curvature and
torsion,
%
\begin{align}
  g^{ab}\,\partial_a\xi\,\partial_b\xi
  \;\approx\;
  &\bigl(1 + 2\kappa x\bigr)\,(\partial_z\xi)^2
  + (\partial_x\xi)^2
  + (\partial_y\xi)^2
  \nonumber\\[0.5ex]
  &\quad
  + 2\tau y\,(\partial_z\xi)(\partial_x\xi)
  - 2\tau x\,(\partial_z\xi)(\partial_y\xi),
  \label{eq:xi-gradient-tubular}
\end{align}
%
so that the gradient energy density of the amplitude mode takes the
schematic form
%
\begin{equation}
  \mathcal{E}_{\mathrm{grad}}
  \;\approx\;
  \frac{T_\mathrm{eff}}{2}\Bigl[
    \bigl(1+2\kappa x\bigr)(\partial_z\xi)^2
    + (\partial_x\xi)^2
    + (\partial_y\xi)^2
    + 2\tau y\,(\partial_z\xi)(\partial_x\xi)
    - 2\tau x\,(\partial_z\xi)(\partial_y\xi)
  \Bigr].
  \label{eq:xi-gradient-energy}
\end{equation}
%
Two features are worth emphasizing. First, the curvature term
$\propto\kappa x(\partial_z\xi)^2$ implies that for sufficiently
large longitudinal gradients of $\xi$ it becomes energetically
favourable for the mode to shift its support away from the
centerline in the transverse directions $(x,y)$, i.e.\ to develop a
lateral ``bulge''. Second, the torsion terms
$\propto\tau\,(\partial_z\xi)(\partial_x\xi)$ and
$\propto\tau\,(\partial_z\xi)(\partial_y\xi)$ show that a purely
longitudinal oscillation is not an eigenmode once the tube is twisted:
steep variations along $z$ inevitably source transverse gradients
$\partial_x\xi$ and $\partial_y\xi$. Together, these geometric
couplings provide a concrete mechanism by which the strongly excited,
Compton-frequency internal mode of the toroidal electron drives lateral
deformations of the brane that confine the energy to a finite tubular
region, while the microscopic spring rest length controls the relative
cost of longitudinal stretching versus transverse bulging through the
effective tension $T_\mathrm{eff}$ and thus sets the critical amplitude
threshold for self-confinement.

\paragraph{Gauge sector from Berry/Wilczek--Zee holonomy.}

In this work, we do \emph{not} close the electromagnetic
sector by introducing a phenomenological Poisson equation for a scalar
potential derived from a coarse-grained transverse-displacement field.
Instead, the $U(1)$ gauge structure is reconstructed from the Berry
connection of band-isolated substrate modes, as detailed in
Section~\ref{conceptual-model}. The internal circulating mode of the
toroidal electron acts as a source of Berry curvature and
holonomy defects; the effective electromagnetic potential $A_\mu^{\text{eff}}$
is identified with this connection. The transverse embedding
deformation (represented by $X^4$ in this gauge) still encodes energy density and contributes to the
induced metric $g_{ij}[\mathbf{X}]$ (gravitational sector), but the
electromagnetic sector is traced to the geometric phase structure of
the polarization bundle, not directly to the Laplacian of a displacement component.


\section{Radial-slice confinement constraints and effective index profile}
\label{app:confinement-slice}

This appendix formalizes a simple ``radial slice'' view of confinement for the toroidal electron.
The goal is not to solve the full nonlinear field equations in closed form, but to identify the
\emph{static background} features that must be present for a photon-like narrowband mode to circulate
around a ring of radius $R$ without rapid leakage.

\paragraph{Geometry vs.\ propagation.}
The embedding $\mathbf X(x,t)\in\mathbb R^4$ (and thus the induced metric $g_{ij}=\partial_i\mathbf X\cdot\partial_j\mathbf X$)
is a geometric object. By contrast, an ``effective refractive index'' profile $n_{\rm eff}(x)$ is a
property of \emph{waves on that background}: it is defined by the dispersion relation obtained from
linearizing the substrate dynamics around a (possibly slowly varying) background configuration
$\mathbf X_0(x)$.
A finite maximal propagation speed of the medium is a kinematic constraint, but by itself does not
produce bending; bending requires spatial gradients of the effective propagation characteristics.

\paragraph{Local frame and radial line.}
Fix a point on the torus centerline and introduce a local orthonormal frame
$\{\hat e_s,\hat e_r,\hat e_\phi,\hat e_4\}$, where $\hat e_s$ is tangent along the torus,
$\hat e_r$ points radially outward in the brane (perpendicular to the centerline),
$\hat e_\phi$ completes an orthonormal basis in the brane cross-section, and $\hat e_4$ denotes the
local embedding-normal direction (the $X^4$ direction in a convenient gauge).
Consider the 1D line parameterized by $r\ge 0$ along $\hat e_r$ starting at the centerline.

\paragraph{Static background ansatz.}
Along this line we write a minimal static deformation in terms of three scalar profiles:
an embedding-normal bulge $w(r)$ (primarily in $X^4$) and two in-brane displacements
$u_r(r)$ (along $\hat e_r$) and $u_\phi(r)$ (along $\hat e_\phi$),
\begin{equation}
\mathbf X_0(r) \;\approx\; \mathbf X_\star(r)
+ u_r(r)\,\hat e_r + u_\phi(r)\,\hat e_\phi + w(r)\,\hat e_4,
\label{eq:radial-slice-embedding}
\end{equation}
where $\mathbf X_\star$ denotes the reference embedding (e.g.\ the pre-tensioned ground state).

The induced metric component along the radial direction is then
\begin{equation}
g_{rr}(r)=|\partial_r\mathbf X_0|^2
\approx (1+u_r'(r))^2 + (u_\phi'(r))^2 + (w'(r))^2 .
\label{eq:radial-slice-grr}
\end{equation}
Along the torus direction $s$ one obtains an additional geometric factor from curvature
$\kappa=1/R$; to leading order, offsets across the tube modify the local arclength by
\begin{equation}
g_{ss}(r)\;\propto\;\bigl(1-\kappa r\bigr)^2,
\label{eq:radial-slice-gss-curvature}
\end{equation}
up to strain-dependent corrections from the full embedding.

\paragraph{Stress and effective wave speed.}
Given a stretch energy density $W(g;g^0)$, define the stress-like tensor
\begin{equation}
P^{ij}(r)=2\,\frac{\partial W}{\partial g_{ij}}(r).
\label{eq:radial-slice-stress}
\end{equation}
For tension-dominated branches propagating primarily along the torus direction, a useful proxy for
the local effective wave speed is
\begin{equation}
c_{\rm eff}^2(r)\;\propto\;\frac{P^{ss}(r)}{\rho_m},
\label{eq:radial-slice-ceff-proxy}
\end{equation}
where $\rho_m$ is the mass density. More generally, $c_{\rm eff}(r)$ is determined by the eigenvalues
of the linearized operator around $\mathbf X_0$; Eq.~\eqref{eq:radial-slice-ceff-proxy} captures the
basic dependence on the background stress. We define an effective refractive index
\begin{equation}
n_{\rm eff}(r)=\frac{c_{\rm ref}}{c_{\rm eff}(r)},
\label{eq:radial-slice-neff}
\end{equation}
for a chosen reference speed $c_{\rm ref}$.

\paragraph{Order-of-magnitude bending requirement.}
Treat the circulating mode as guided in a tube of transverse radius $a$ that is bent into a ring of
radius $R$. A robust waveguide estimate is that bending without rapid leakage requires an index
contrast across the guided region of order
\begin{equation}
\frac{\Delta n}{n}\;\sim\;\frac{a}{R},
\label{eq:radial-slice-bend-criterion}
\end{equation}
equivalently $\Delta c/c\sim a/R$. This condition expresses that curvature must be supported by a
transverse gradient of the effective propagation characteristics (self-induced waveguiding).

\paragraph{Inner/outer coherence for a double ridge.}
If the electron core supports two intertwined ridges (``double winding'') separated by a transverse
offset $\pm w_0$ from the tube center, the inner and outer paths have different lengths
$2\pi(R\mp w_0)$. For a coherent closed mode at a common temporal frequency $\omega$, closure requires
different local phase advance per arclength,
\begin{equation}
k_{\rm in}(R-w_0)\;\approx\;k_{\rm out}(R+w_0),
\qquad
\frac{\Delta k}{k}\;\sim\;\frac{2w_0}{R}.
\label{eq:radial-slice-inner-outer-k}
\end{equation}
Thus it is typically \emph{not} the temporal frequency that varies across the cross-section, but the
local wavenumber (or phase velocity) at fixed $\omega$.
The same background stress profile that guides the mode must therefore also provide a cross-sectional
dispersion gradient compatible with Eq.~\eqref{eq:radial-slice-inner-outer-k}.

Equations~\eqref{eq:radial-slice-bend-criterion}--\eqref{eq:radial-slice-inner-outer-k} provide
direct, testable constraints on the deformation profiles $w(r)$, $u_r(r)$, $u_\phi(r)$ through their
effect on $P^{ss}(r)$ and hence on $n_{\rm eff}(r)$.
